\section{The triplet is modular}\label{sec:modular}

\paragraph{The question.}
The triplet of Section~\ref{sec:weave} is acted on by the moves. The next question is what that action is. If it were
an arbitrary finite group, the triplet would carry no further structure. We find instead that the action is the
multiplier of a classical family of theta functions: the grammar's moves, with no modular input, reproduce the level-four
theta multiplier that modular flavour models use, up to a twist by a power of the eta character. The twist doubles the
image from order $48$ to order $96$ and moves the allowed weights to $3/2$ modulo $2$ (Theorem~\ref{thm:theta}). The record is cited by arc~\cite{OriginAxiomRepo}, and what is known is
cited where it is used.

\paragraph{The setting.}
Take the moves on the ruled branch, that is $\langle L,R\rangle$ with no swap and no sign move beyond those generated
(Section~\ref{sec:grammar}). This group is the braid group $B_3$, acting on the once-punctured torus group $F_2=\langle
a,b\rangle$, and it maps onto $\SL(2,\Z)$ with kernel $\langle\Delta^4\rangle$, where $\Delta^{-1}=RL^{-1}R$ and
$\Delta^2\mapsto-I$ on homology. The metaplectic group is $\mathrm{Mp}_2(\Z)=\langle s,t\mid s^2=(st)^3,\ s^8=1\rangle=B_3/\langle\Delta^8\rangle$.
Write $s$ and $t$ for the images of $RL^{-1}R$ and $L$ in $B_3/\langle\Delta^8\rangle=\mathrm{Mp}_2(\Z)$.

\begin{proposition}[the inner automorphisms are not moves, B1634]\label{prop:inner}
The group $\langle L,R\rangle\cong B_3$ meets $\mathrm{Inn}(F_2)$ exactly in $\langle\Delta^4\rangle=\langle\mathrm{conj}(c)\rangle$, $c=a\,b^{-1}a^{-1}b$.
To word length $12$ the words acting on $F_2$ as inner automorphisms are $1{,}944$ relators (the identity), $196$ words for
$\mathrm{conj}(c)$ and $196$ for its inverse, and none for $\mathrm{conj}(a^{\pm1})$ or
$\mathrm{conj}(b^{\pm1})$. Every word acting on $F_2$ as an inner automorphism lifts to a scalar on the triplet.
\end{proposition}
\noindent\emph{Proof and scope.} The intersection is the classical kernel of $B_3\to\SL(2,\Z)$ acting on $\pi_1$ of the
punctured torus, and the census to length $12$ is exact (B1634, a verification of a claim made within the program). The
statement is about the ruled branch only. On branches that add the sign move or the swap, or give the records' base
point moves of its own, the inner automorphisms act, and results that use them are stated there.

\paragraph{The multiplier.}
Let $T$ be the triplet of Theorem~\ref{thm:triplet}, with the lifts of the moves defined on $V=T\oplus\bar T$ by the
parity-twisted action of Section~\ref{sec:weave}. We write $\rho$ for the representation of the moves on $T$,
$\rho(\gamma)=\mathrm{lift}(\gamma)^{-1}$; this is the convention in which $T$ is cohomology, so its monodromy is
the pull-back's inverse. The theta constants are $\theta_2,\theta_3,\theta_4$ in the nome $q=e^{i\pi\tau}$, and
$u=(\theta_2^2,\theta_3^2,\theta_4^2)$ satisfies Jacobi's identity $\theta_3^4=\theta_2^4+\theta_4^4$. Its components are
modular forms of weight one on $\Gamma(4)$, and the graded ring of the modular forms of $\Gamma(4)$ is generated by them
modulo that one relation. Both sides have dimension $2d+1$ in weight $d$: $\Gamma(4)$ has genus zero and six regular
cusps~\cite[ch.~3]{DiamondShurman2005}, and the quotient of the polynomial ring by the irreducible quadric injects
because the image has Krull dimension two. Modular flavour models use this ring as their
input~\cite{NovichkovPenedoPetcov2021,LiuYaoDing2020}. The transformation laws of the theta constants,
the theta functions with characteristics on the torus, are classical~\cite{WhittakerWatson1927,AlvarezGaumeMooreVafa1986}.
We write $\rho_u$ for the multiplier of the vector $u$ under $\mathrm{Mp}_2(\Z)$; it is trivial on $\Gamma(4)$, and
$\chi_\eta$ for the multiplier character of the Dedekind eta function.

\begin{theorem}[the theta multiplier, T3]\label{thm:theta}
Take the ruled weave on the triplet $T$ in the convention above.
\begin{enumerate}
\item The lifts of the moves satisfy the braid relation and send $\Delta^4$ to $-I$. They are therefore a genuine
  (not projective) representation $\rho$ of $B_3/\langle\Delta^8\rangle=\mathrm{Mp}_2(\Z)$.
\item $\rho$ equals, with no rescaling, the multiplier of the vector $\eta^{21}\cdot(\theta_2^2,\theta_3^2,\theta_4^2)$:
  \[
    \rho=\chi_\eta^{21}\otimes\rho_u ,
  \]
  the weight-one multiplier $\rho_u$ of the $\Gamma(4)$ triplet twisted by a power of the eta character. The intertwiner
  from the parity lines to the theta constants is monomial: $\theta_3^2$ lies on the clock's line (the parity
  character $(-,-)$), $\theta_2^2$ on the line $(+,-)$, and $\theta_4^2$ on the line $(-,+)$. (A line is labelled by the
  signs of its character on $a$ and $b$. Lemma~\ref{lem:parity}, Appendix~\ref{app:hand}'s row-vector right action and
  the record's pull-back of characters share these labels: $L$ fixes the parity $(0,1)$, the line $(+,-)$, and $R$ fixes
  $(1,0)$, the line $(-,+)$. The labels $(+,-)$ and $(-,+)$ exchange only under left multiplication of column vectors.)
  The exponents of $t$ on $T$ are
  $(\tfrac18,\tfrac38,\tfrac78)$ turns.
\item The dimension of the space $M_k(\rho)$ of vector-valued forms of weight $k$ for $\rho$ is
  $0,0,1,1,2,2,3,3$ at the allowed weights $k=-\tfrac12,\tfrac32,\tfrac72,\dots,\tfrac{27}{2}$ (step $2$; a weight is
  \emph{allowed} when the central element acts consistently with it, a necessary condition, not the existence of a
  form). In particular $\dim M_{3/2}(\rho)=0$: weight $\tfrac32$ is allowed and carries no form. The dual
  representation is not allowed at weight $\tfrac32$.
\end{enumerate}
\end{theorem}
\begin{proof}[Proof sketch (record B1634, with B1640's convention scoping)]
(1) The spin lifts of the moves are elements of the binary octahedral group $2O\subset\mathrm{SU}(2)$ (B1634). Acting on
$V=T\oplus\bar T$ through the parity twist they generate a finite group of order $96$, which acts faithfully on $T$
because $T$ and $\bar T$ have the same kernel (B1600). The braid relation holds exactly: $g_Lg_R^{-1}g_L=g_R^{-1}g_Lg_R^{-1}=(i+j)/\sqrt2$ in $2O$ for the spin lifts
$g_L=(1+i)/\sqrt2$ and $g_R=(1-j)/\sqrt2$, and $LR^{-1}L=R^{-1}LR^{-1}$ in $\mathrm{Aut}(F_2)$ (the record also checks it
numerically, to $3\cdot10^{-16}$); and $\Delta^4$, the generator of the kernel, acts as $-1$ (turns $\tfrac18,\tfrac18,\tfrac58$ for the
lift of $\Delta$). The lifts compose in reverse, so $\rho=\mathrm{lift}^{-1}$ is a homomorphism. (2) With $s,t$ as above,
$\rho(s)$ and $\rho(t)$ satisfy $s^2=(st)^3$ and $s^8=1$ exactly, so $\rho$ is a representation of $\mathrm{Mp}_2(\Z)$. The
theta transformation laws, evaluated from the series in $\tau$ at several points, fix $\rho_u$, and a linear solve for an
intertwiner with $\rho$ returns an intertwiner with no rescaling (both scale factors equal to $1$) and smallest singular value $1.5\cdot10^{-15}$; its monomial shape is
exact. Four of the eight identifications of $(s,t)$ with the generators and their inverses need no rescaling, and two of
them are allowed at weight $3/2$. The record's sealed check of the classical laws failed at one of its two points (the
$T$ law off by a relative $1.10$ at $\tau=0.13+1.21i$); it was traced to the instrument's principal branch of $q^{1/4}$,
and at the other point both laws hold to $2\cdot10^{-15}$ (B1634, disclosed there). (3) If $f\in M_k(\rho)$ then $\eta^3f$ is a form of weight $k+\tfrac32$ for the
multiplier of $u$ vanishing at the cusp, and the converse holds; the forms of $\Gamma(4)$ are the polynomials in $u$
modulo the one relation ($2d+1$ of them in degree $d$), and the count is exact linear algebra. At $k=\tfrac32$ the only
candidate is the Serre derivative of $u$, whose constant term $-\tfrac1{12}(0,1,1)$ does not vanish, so it gives no form
and $\dim M_{3/2}(\rho)=0$. As a cross-check, the counts above weight $2$ agree with the dimension given by the
Selberg trace formula~\cite[\S7]{Borcherds2000}. Below weight $2$ the same expression is only an Euler characteristic,
$\dim M_k(\rho)-\dim S_{2-k}(\rho^\vee)$; at $k=\tfrac32$ it is $0$, and both terms vanish.
\end{proof}

\noindent\emph{Scope.} Four limitations are part of the statement.
\begin{itemize}
\item \emph{What is classical, what is automatic and what is computed.} The modular group permutes the three even theta
  characteristics of the torus through $\SL(2,\mathbb F_2)\cong S_3$, and the theta multiplier of Section~\ref{sec:modular}'s
  setting is classical. Two genuine representations of $\mathrm{Mp}_2(\Z)$ that agree projectively differ by a character,
  and the characters of $\mathrm{Mp}_2(\Z)$ are the powers of $\chi_\eta$ (its abelianisation is $\Z/24$). So once the
  parity lines are identified with the even characteristics, a power of $\chi_\eta$ is automatic. What the theorem
  computes is the projective agreement (the relative phases of the monomial matrices), the exponent $21$, and with it the
  central character that allows weight $\tfrac32$; it computes a stated lift, not a derived flavour symmetry.
\item \emph{The hand is a convention.} Writing $\rho=\mathrm{lift}^{-1}$ is one of two homomorphisms; the other is the
  complex conjugate, and the choice of which piece of $V$ is called $T$ is a second binary choice. The two together
  select $\rho$ or $\bar\rho$. The representation is irreducible and not equivalent to its conjugate (an independent exact
  audit, B1640). The record's choice has weight $3/2$ allowed (in the sense above; no form exists there) and admits the lowest
  symmetric weight-one coupling (B1640); the conjugate has neither. So the index and that coupling come with the same hand,
  and which hand the matter carries is not derived here. The hand's functorial reason (cohomology pulls back) is a
  reading of the record, not a theorem of this section.
\item \emph{The eta twist.} The theorem identifies $\rho$ with the theta multiplier up to the one-dimensional character
  $\chi_\eta^{21}$; it is not the bare weight-one multiplier of $\Gamma(4)$'s triplet. The restriction of $\rho$ to $\Gamma(4)$ is that
  character. We do not claim that the twist is forced rather than a convention of the homology generators.
\item \emph{The dimension formula.} The identification of the forms of $\Gamma(4)$ with the polynomials in $u$ modulo
  Jacobi's identity holds in every degree $d$ by the argument in the setting above (injectivity, and dimension $2d+1$ on both
  sides). It rests on the standard dimension formula $2d+1$ for the genus-zero curve with six regular cusps; it was also
  certified numerically for $d\le5$.
\end{itemize}

\begin{remark}[the modulus breaks the weave's group]\label{rem:modulus}
On the ruled branch the weave forces no point of the upper half-plane. At a generic $\tau$ the stabilizer of $\tau$ in
$\langle L,R\rangle$ is the centre, so the residual symmetry on $T$ is four scalars. At $\tau=i$ it is cyclic of order $8$,
a plane containing the clock's line and a line in the span of the letters' two lines; at $\tau=\omega=e^{2\pi i/3}$ it is cyclic of order $12$, with
three eigenlines, each with squared overlap $\tfrac13$ with every parity line (B1634). Consequently the modulus $\tau$ is the state that
breaks the weave's group. On the modular-flavour ansatz (holomorphic couplings, a minimal K\"ahler term, one copy of $T$
per field, one Higgs per sector), which this paper uses only for the ruled-branch bound of Theorem~\ref{bd:reducesnone},
the flavour numbers would be the values of vector-valued modular forms for $\rho$ at $\tau$. The couplings of
Section~\ref{sec:couplings} are built from such forms together with norms that are not holomorphic, so they are not of that
kind. The
statements of the earlier record that used the inner automorphisms as symmetries (an order-48 residual at $\omega$, a
permutation mixing for every $\tau$) hold on the branches where those act and not on the ruled branch.
\end{remark}

\subsection{The lowest couplings, the determinant character and the strong angle}\label{sec:anomaly}

This subsection works on the ruled branch with the modular-flavour ansatz of Remark~\ref{rem:modulus}, with one singlet
Higgs of weight zero, so that a sector's coupling between fields carried by $\rho_L$ and $\rho_R$ is a form in
$M_k(\rho_L^\vee\otimes\rho_R^\vee)$: odd $k$ when both fields carry $T$, even $k$ when one carries $\bar T$. The
statements about the strong angle add one input, the moves as a symmetry of the fermions (Table~\ref{tab:inputs}).

\begin{corollary}[the half floor at weight one, B1635, B1636]\label{cor:halffloor}
With every field in $T$, the coupling of weight one is unique up to a constant (B1635) and is the theta constants placed
off the diagonal,
\[
  Y_1\propto\begin{pmatrix}0&\theta_4^2&\theta_3^2\\ \theta_4^2&0&\theta_2^2\\ \theta_3^2&\theta_2^2&0\end{pmatrix}.
\]
Its singular values $s_1\ge s_2\ge s_3$ satisfy $s_1=s_2+s_3$ at every $\tau$, so $s_2/s_1\ge\tfrac12$ at every
$\tau$. The spectrum is $(1,\tfrac12,\tfrac12)$ at $\omega$ and $(1,\varphi^{-1},\varphi^{-2})$ at $i$. With one form per
sector, $M_u\propto M_d$ and the mixing matrix is the identity.
\end{corollary}
\begin{proof}
For a complex symmetric matrix with zero diagonal, the elementary symmetric functions of the eigenvalues of $YY^\dagger$
satisfy $e_1^2=4e_2$, a polynomial identity in the three entries (checked symbolically). In terms of the singular values
this is $(s_1+s_2+s_3)(-s_1+s_2+s_3)(s_1-s_2+s_3)(s_1+s_2-s_3)=0$, so the largest is the sum of the other two, and
$s_2\ge s_3$ gives $s_2\ge s_1/2$. The coupling's form and the spectrum at $i$ and on the tick's geodesic are B1636's,
and the sum rule for every coupling of that piece is its census G2. At $\omega$ the three theta constants have equal
moduli, since $\lambda(\omega)=e^{i\pi/3}$ gives $|\theta_2^4|=|\theta_4^4|=|\theta_3^4|$; the spectrum computed there
is $(1,\tfrac12,\tfrac12)$.
\end{proof}
\noindent\emph{Scope.} The floor is a statement about the ruled branch's parameter-free coupling under the ansatz; the
data need $m_2/m_1$ below $0.06$ in every charged sector. With $R$ in $\bar T$ the coupling of weight zero is the
identity, three degenerate masses (B1635). In the fibre frame the floor $\tfrac12$ of Theorem~\ref{thm:isotropy}(2)
holds for real Higgs directions only; complex directions reach any ratio near the isotropic cone. So the floor is not a
property of every frame.

\paragraph{The determinant character.}
\begin{proposition}[the determinant character, B1674]\label{prop:detchar}
$\det\rho=\chi_\eta^{9}$.
\end{proposition}
\begin{proof}
A character of $\mathrm{Mp}_2(\Z)$ is fixed by its value on $t$, whose image generates the abelianisation $\Z/24$, and
$\chi_\eta(t)=e^{2\pi i/24}$. By Theorem~\ref{thm:theta}(2) the exponents of $t$ on $T$ sum to $\tfrac18+\tfrac38+\tfrac78
=\tfrac{11}8\equiv\tfrac{9}{24}$ turns. The value on $s$ agrees (B1674).
\end{proof}

\paragraph{The moves' anomaly.}
Suppose the moves act on the Standard Model's fermions as a symmetry. Each multiplet $f\in\{Q,u^c,d^c,L,e^c\}$, and
$\nu^c$ where it is named, carries its three generations in $T$ ($s_f=+1$) or in $\bar T$ ($s_f=-1$). A modular
transformation then rotates the left-handed fermions, and the rotation shifts the angle of each gauge factor $G$ by its
anomaly. At the level of the character the shift is the character $\chi_\eta^{n_G}$, with
\[
  n_G(s)=9\sum_f w_G(f)\,s_f\pmod{24},
\]
where $w_G(f)$ is twice the Dynkin index of $f$'s representation of $G$ times its multiplicity under the other factors:
$(w_Q,w_{u},w_{d})=(2,1,1)$ for $\mathrm{SU}(3)^2$, $(w_Q,w_L)=(3,1)$ for $\mathrm{SU}(2)^2$, and the dimensions
$(6,3,3,2,1,1)$ for gravity.

\begin{proposition}[the anomaly of the moves, B1674]\label{prop:anomaly}
Given the triplet postulate and the moves as a symmetry of the fermions, at the level of the character:
\begin{enumerate}
\item the mixed anomaly with $\mathrm{SU}(3)$ vanishes exactly when $s_u=s_d=-s_Q$, that is, when the moves act on the
  left-handed and the right-handed quarks by the same representation (16 of the 64 assignments);
\item the mixed anomaly with $\mathrm{SU}(2)$ never vanishes: $n_{\mathrm{SU}(2)}\in\{6,12,18\}$;
\item among the assignments of (1), the gravitational anomaly vanishes exactly when $s_e=s_\nu=-s_L$, and without
  $\nu^c$ it never vanishes.
\end{enumerate}
\end{proposition}
\begin{proof}
Since $3$ is invertible modulo $8$, $n_G\equiv0\pmod{24}$ exactly when $\sum_f w_G(f)s_f\equiv0\pmod 8$. In (1) the sum
$2s_Q+s_u+s_d$ lies in $\{0,\pm2,\pm4\}$ and vanishes exactly as stated. In (2) the sum $3s_Q+s_L$ lies in
$\{\pm2,\pm4\}$. In (3), given (1), the sum is $2s_L+s_e+s_\nu$, which vanishes exactly when $s_e=s_\nu=-s_L$; without
$\nu^c$ it is odd. B1674 enumerates the 64 assignments exactly.
\end{proof}

\paragraph{The couplings' determinants and the strong angle.}
The strong angle is $\bar\theta=\theta-\arg\det(M_uM_d)$, with $M=Yv$. With one singlet Higgs its phase cancels from
$\bar\theta$, and a move shifts $\bar\theta$ by the phase of the multiplier of $\det(Y_uY_d)$. That is the anomaly of
Proposition~\ref{prop:anomaly} together with the cocycle of the weights.

\begin{proposition}[the determinants and the strong angle, B1674]\label{prop:strong}
Let the up and down couplings have weights $k_u,k_d$, and put $K=k_u+k_d$.
\begin{enumerate}
\item $\det(Y_uY_d)$ is a form of weight $3K$ with character $\chi_\eta^{m}$, $m=-9(2s_Q+s_u+s_d)$. It is $\eta^{6K}$ times
  a modular function exactly when $6K\equiv m\pmod{24}$:
  \begin{itemize}
  \item for every quark in $T$, exactly when $K\equiv2\pmod 4$;
  \item for $s_u=s_d=-s_Q$, exactly when $K\equiv0\pmod 4$;
  \item for the mixed assignments, when $K\equiv1$ or $3\pmod4$, by the hand.
  \end{itemize}
  When it holds, a counterterm $\theta(\tau)=\theta_0+6K\arg\eta(\tau)$, which we call the Dedekind counterterm, makes
  $\bar\theta$ invariant under the moves, and $\bar\theta-\theta_0=-\arg g$ with $g=\det(Y_uY_d)/\eta^{6K}$.
\item At weight one with every quark in $T$, $\det Y_1=2\theta_2^2\theta_3^2\theta_4^2=8\eta^6$. With real coefficients
  and no counterterm, $\bar\theta-\theta\equiv-12\arg\eta(\tau)\pmod\pi$. This is $-\pi\,\mathrm{Re}\,\tau$ up to terms
  of order $e^{-2\pi\,\mathrm{Im}\,\tau}$, and it is $\pi/2$ exactly on $\mathrm{Re}\,\tau=\tfrac12$. With the counterterm,
  $g$ is constant, so $\bar\theta$ takes one value at every $\tau$.
\item At weight three with every quark in $T$, and at weight two with $s_u=s_d=-s_Q$, each sector's determinant is a
  constant times $\eta^6E_6$, and a constant times $E_6$, respectively.
  \begin{itemize}
  \item With the counterterm, $g$ is a constant times $E_6^2/\Delta=j-1728$. With real coefficients, $\bar\theta-\theta_0$
    is a multiple of $\pi$ exactly on the imaginary axis and on the boundary of the fundamental domain. These are the
    loci where generalised CP is conserved in modular models~\cite{NovichkovPenedoPetcovTitov2019}, and there the mixing
    phase vanishes.
  \item Without the counterterm, at weight three, $\bar\theta-\theta\equiv-(12\arg\eta+2\arg E_6)\pmod\pi$. This vanishes
    on the imaginary axis and also on a curve inside the fundamental domain: from $\mathrm{Re}\,\tau\approx0.324$ at
    $\mathrm{Im}\,\tau=1$ to $\mathrm{Re}\,\tau\approx0.172$ at $\mathrm{Im}\,\tau=1.2$.
  \end{itemize}
\end{enumerate}
\end{proposition}
\begin{proof}[Proof sketch]
(1) The coupling transforms as $Y(\gamma\tau)=(c\tau+d)^k\rho_L^\vee(\gamma)Y(\tau)\rho_R^\vee(\gamma)^{T}$, so its
determinant carries $(c\tau+d)^{3k}(\det\rho_L\det\rho_R)^{-1}$, and Proposition~\ref{prop:detchar} gives $m$. The power
$\eta^{6K}$ carries $\chi_\eta^{6K}(c\tau+d)^{3K}$, so the quotient has weight zero and trivial character exactly when
$6K\equiv m$.

(2) The determinant of a zero-diagonal symmetric matrix is twice the product of its entries, and
$\theta_2\theta_3\theta_4=2\eta^3$~\cite{WhittakerWatson1927}; the identity and the phase law are also checked
independently in B1671. On $\mathrm{Re}\,\tau=\tfrac12$ the nome
$e^{2\pi i\tau}$ is negative, so $\eta^{12}$ is $i$ times a positive number.

(3) The quotient of the determinant by $\eta^6$, or the determinant itself, is a form of weight six for $\SL(2,\Z)$:
its exponents at the cusp are integers no smaller than $-\tfrac14$, hence non-negative. The space of such forms is
spanned by $E_6$ alone. Since $E_4^3-E_6^2=1728\Delta$, the quotient $g$ with the counterterm is a constant times
$E_6^2/\Delta=j-1728$, and $j$ is real on the fundamental domain exactly on the imaginary axis and on its
boundary~\cite{DiamondShurman2005}. B1674 checks every identity at five moduli to $10^{-28}$, and the zero sets on a
grid. The interior curve was first seen while reading the record, before B1674's rule, and B1674 restates it rather than
predicting it.
\end{proof}
\noindent\emph{Scope.} Several limitations are part of the statements.
\begin{itemize}
\item The moves as a symmetry of the fermions and $\tau$ as a field are inputs, and so is the Dedekind counterterm
  (Table~\ref{tab:inputs}). In string compactifications a term of this form comes with the threshold corrections that
  cancel the target-space modular anomaly~\cite{DixonKaplunovskyLouis1991,IbanezLust1992}.
\item The strong angle $\bar\theta$ is one of the nineteen, and it is not fixed here. With the counterterm it is the
  constant $\theta_0$, a free input. Without it, with real coefficients, it is $\theta-\pi\,\mathrm{Re}\,\tau$ up to
  terms of order $e^{-2\pi\,\mathrm{Im}\,\tau}$ (B1671), tied to $\mathrm{Re}\,\tau$ but not fixed by it. The count
  stays at 0 of 19.
\item The strong-angle statements are tree-level, on the holomorphic ansatz with real coefficients and the read-out of
  Table~\ref{tab:inputs}. The mixing phase on the interior curve of (3) is not computed.
\item $\mathrm{SU}(2)$'s angle is not physical in the Standard Model, so Proposition~\ref{prop:anomaly}(2) binds only if
  the moves are gauged. The condition with $\mathrm{U}(1)_Y$ depends on the normalisation and is not stated.
\item One control of B1674 failed as worded: the sign of the determinant character on words whose image has
  $|\mathrm{Re}\,\gamma\tau|>1$, through the instrument's principal branch of $q^{1/4}$, as in B1634. It was recomputed
  without a branch after the run (Appendix~\ref{app:corrections}).
\end{itemize}
The known escape from the obstruction in (3) keeps determinants of weight zero, with reducible representations and
unequal weights~\cite{FeruglioParriciatuStrumiaTitov2024,PenedoPetcov2024}. With the irreducible triplet it is not
available, and Proposition~\ref{prop:strong} states what is: at weight one the angle is constant, and the masses are
excluded by Corollary~\ref{cor:halffloor}; at weights two and three, with real coefficients, $\bar\theta-\theta_0$ is a multiple of $\pi$ only where
generalised CP is conserved.

\paragraph{Position against the literature.}
We state which parts are known and which are the paper's.
\begin{itemize}
\item \emph{Known.} Modular flavour symmetry, in which fermion mass matrices are modular forms and the flavour group
  is a finite quotient of the modular group acting on a multiplet~\cite{Feruglio2019}. The weight-one level-four forms
  from two theta constants, and the triplet of $S_4'$ they form, as the inputs of flavour models~\cite{NovichkovPenedoPetcov2021,LiuYaoDing2020}.
  Residual symmetries at the fixed points of the modular group and the hierarchies that follow from the neighbourhood
  of a fixed point~\cite{NovichkovPenedoPetcov2021hier}. Zero modes of magnetized tori as multiplets of a
  modular group, and Yukawa couplings from their overlaps~\cite{CremadesIbanezMarchesano2004,KobayashiTamba2019}. The
  theta constants as the three even theta characteristics of the torus~\cite{AlvarezGaumeMooreVafa1986}.
\item \emph{The paper's.} Not the triplet, nor the modular group, nor the theta constants. The paper's content is the
  chain: the records' own parities and the common point give the triplet (Theorem~\ref{thm:triplet}), the
  grammar's moves act on it as a genuine $\mathrm{Mp}_2(\Z)$ representation (Theorem~\ref{thm:theta}), and that
  representation is the theta multiplier up to a monomial intertwiner and an eta character. Modular flavour models
  choose the untwisted symmetry; here the symmetry of the moves is that one up to the eta twist and the hand, as scoped above.
\end{itemize}
The theorem does not say that the program's weave is any of those models, and it does not select a modulus, a weight
or a vacuum. The reading of the zero modes as the overlap integrals of the magnetized-torus literature is an open lead of the
record, not a result of this section.

\reading{Three objects permuted by an order-three element, with a metaplectic group acting on them, appear on both
sides: the records' parities and the theta constants. We read the triplet as modular because the records are:
the symmetry is the one the moves carry, and a flavour model that chooses it is choosing, up to the stated twist (which moves the weights to $3/2$ modulo $2$) and
hand, what the grammar already does. What this reading cannot supply is the modulus. A symmetry that the moves carry
is broken by the state, and which state is a selection (Section~\ref{sec:selection}).}
